% !TEX TS-program = pdflatex
% !TeX program = pdflatex
% !TEX encoding = UTF-8
% !TEX spellcheck = fr

\documentclass[xcolor=table, usenames,dvipsnames]{beamer}


%\usepackage{fullpage}
%\usepackage[left=2.8cm,right=2.2cm,top=2 cm,bottom=2 cm]{geometry}
\setbeamersize{text margin left=10pt,text margin right=10pt}
\usepackage{amsmath,amssymb} 
\usepackage[T1]{fontenc}
\usepackage[utf8]{inputenc}
\usepackage[english,french]{babel}
\usepackage{txfonts}
\usepackage[]{graphicx}
\usepackage{multirow}
\usepackage{hyperref}
\usepackage{colortbl}
\usepackage{listings}
\usepackage{wrapfig}
\usepackage{multicol}
\usepackage[export]{adjustbox} %for images in table, also for frame

\hypersetup{
	colorlinks,
	urlcolor = blue
}

%\renewcommand{\baselinestretch}{1.5}

\def\supit#1{\raisebox{0.8ex}{\small\it #1}\hspace{0.05em}}

\AtBeginSection{%
	\begin{frame}
		\sectionpage
	\end{frame}
}

\newcommand{\rottext}[2]{%
	\rotatebox{90}{%
	\begin{minipage}{#1}%
		\raggedleft#2%
	\end{minipage}%
	}%
}

\usepackage{longtable}
\usepackage{tabu}


\institute{ %
École  nationale Supérieure d'Informatique (ESI, ex. INI), Algérie
}
\author[ \textbf{\footnotesize  \insertframenumber/\inserttotalframenumber} \hspace*{\fill} ESI (2019-2020)] %
{ARIES Abdelkrime}
%\titlegraphic{\includegraphics[height=1cm]{../img/esi-logo.png}%\hspace*{4.75cm}~


\date{Année unniversitaire: 2019/2020} %\today

\usetheme{Warsaw} % Antibes Boadilla Warsaw

\beamertemplatenavigationsymbolsempty

%\setbeamertemplate{headline}{}

\definecolor{lightblue}{HTML}{D0D2FF}
\definecolor{lightyellow}{HTML}{FFFFAA}
\definecolor{darkblue}{HTML}{0000BB}
\definecolor{olivegreen}{HTML}{006600}
\definecolor{violet}{HTML}{6600CC}

\newcommand{\keyword}[1]{\textcolor{red}{\bfseries\itshape #1}}
\newcommand{\expword}[1]{\textcolor{olivegreen}{#1}}
\newcommand{\optword}[1]{\textcolor{violet}{\bfseries #1}}

\makeatletter
\newcommand\mysphere{%
	\parbox[t]{10pt}{\raisebox{0.2pt}{\beamer@usesphere{item projected}{bigsphere}}}}
\makeatother

%\let\oldtabular\tabular
%\let\endoldtabular\endtabular
%\renewenvironment{tabular}{\rowcolors{2}{white}{lightblue}\oldtabular\rowcolor{blue}}{\endoldtabular}


\NoAutoSpacing %french autospacing after ":"

\def\graphpath{}

\newcommand{\changegraphpath}[1]{\def\graphpath{#1}}

\newcommand{\graphintable}[2]{%
	\includegraphics[height=#1,valign=t]{\graphpath #2}%
}

\newcommand{\vgraphpage}[2][.85\textheight]{%
	\begin{center}%
		\includegraphics[height=#1]{\graphpath #2}%
	\end{center}%
}

\newcommand{\hgraphpage}[2][\textwidth]{%
%	\begin{center}%
		\includegraphics[width=#1]{\graphpath #2}%
%	\end{center}%
}

\newcommand{\insertbibliography}[2]{
	\appendix
	\section*{Bibliographie}
	\nocite{#2}
	
	\begin{multicols*}{2}[\frametitle{\insertsection} \usebeamertemplate{frametitle}]%\usebeamertemplate*{frametitle}\frametitle{Références}
		\tiny
		\bibliography{#1}
		\bibliographystyle{acm}
	\end{multicols*}
}


\usepackage{calligra}

\title[BWEB : 05- Présentation] %
{Bureautique et Web \\Chapitre 05 : Présentation}  

\changegraphpath{../img/presentation/}

\begin{document}

\begin{frame}

\frametitle{Présentation}
\framesubtitle{Pourquoi faire une bonne présentation ?}
	
	\begin{itemize}
		\item Communiquer ses idées : inutile de faire un bon travail sans que les autres le comprennent
		\item Se présenter : avoir des bonnes impressions 
		\item Discuter son travail : éduquer, inspirer, motiver et avoir du retour d'information (feedback)
		\item Mieux comprendre son travail : Lorsqu'on présente son travail, on doit transmettre plus d'informations d'une façon simple (Apprendre par enseignement)
	\end{itemize}
\end{frame}

\begin{frame}
\frametitle{Présentation}
\framesubtitle{Plan}

\begin{multicols}{2}
	%	\small
	\tableofcontents
\end{multicols}
\end{frame}

%===================================================================================
\section{Préparation d'une présentation}
%===================================================================================

\subsection{Objectifs de la présentation}

\begin{frame}
\frametitle{Préparation}
\framesubtitle{Objectif (Pourquoi?)}

\hgraphpage{objectif.pdf}

\end{frame}

\begin{frame}
\frametitle{Préparation Objectif}
\framesubtitle{Information communiquée}

\begin{itemize}
	\item \optword{Présentation informative}
	\begin{itemize}
		\item fournir aux gens des informations sur un concept ou une idée 
		\item exemple : \expword{cours}, \expword{tutoriel}
	\end{itemize}

	\item \optword{Présentation persuasive}
	\begin{itemize}
		\item influencer un changement dans la croyance, l'attitude ou le comportement d'une autre personne 
		\item exemple : \expword{recherche}, \expword{produit}
	\end{itemize}
\end{itemize}

\end{frame}

\begin{frame}
\frametitle{Préparation Objectif}
\framesubtitle{Sens de communication}

\begin{itemize}
	\item \optword{Présentation linéaire}
	\begin{itemize}
		\item seul le présentateur parle lors de la présentation 
		\item but : raconter des idées
		\item exemple, \expword{PFE}, \expword{recherche}
	\end{itemize}

	\item \optword{Présentation interactive}
	\begin{itemize}
		\item l'audience participe lors de la présentation
		\item but : découvrir les idées
		\item exemple, \expword{certains cours}, \expword{démonstration d'un produit}
	\end{itemize}
\end{itemize}

\end{frame}

\subsection{Audience de la présentation}

\begin{frame}
\frametitle{Préparation}
\framesubtitle{Audience (Qui?)}

Il faut poser des questions sur les spectateurs : 
\begin{itemize}
	\item Qui sont-ils?
	\begin{itemize}
		\item Qu'est-ce qui offense l'audience? Comprendre la culture et le contexte dans lesquels l'audience opère
		\item Qu'est-ce qui fait vibrer l'audience? Nouvelles idées, plus d'exemples, illustrations, etc.
		\item Qu'est-ce qui ennuie l'audience? Raconter l'évidence, non maitrise du sujet, répétition, etc.
	\end{itemize}
	\item Que savent-ils?
	\begin{itemize}
		\item Est-ce qu'il y a des experts dans le domaine?
		\item Est-ce qu'ils ont déjà vu une présentation similaire?
	\end{itemize}
	\item Que veulent-ils?
	\begin{itemize}
		\item Plus d'informations ou l'essentiel
		\item Technique ou théorie
	\end{itemize}
\end{itemize}

\end{frame}

\begin{frame}
\frametitle{Préparation : Audience}
\framesubtitle{Types de présentations selon l'audience}

\hgraphpage{audience.pdf}
\end{frame}


\begin{frame}
\frametitle{Préparation : Audience}
\framesubtitle{Expérience}

\begin{itemize}
	\item \optword{Présentation élémentaire}
	\begin{itemize}
		\item audience : non familière avec le thème
		\item exemple : \expword{présenter un langage de programmation à des débutants}
		\item Prévoir des définitions 
		\item Inclure plus d'informations sur le thème
	\end{itemize}
	\item \optword{Présentation approfondie}
	\begin{itemize}
		\item audience : ayant une expérience sur le thème
		\item exemple : \expword{présenter un projet devant un jury des experts}
		\item Moins de définitions 
		\item Moins de fond sur le thème (état de l'art)
		\item Présenter le sujet d'un façon plus profonde (mettre l'accent sur votre travail plus que les travaux similaires)
	\end{itemize}
\end{itemize}

\end{frame}

\begin{frame}
\frametitle{Préparation : Audience}
\framesubtitle{Expérience (Exemple)}

Prenons l'exemple de \cite{carlton-jacob}. Etant donnée la figure suivante, on compare les différents titres et le niveau de compréhension. 

\begin{minipage}{0.60\textwidth}
	\begin{itemize}
		\item \expword{La technologie de synthèse d'ADN s'améliore
			de façon exponentielle} : Sauf les intellos qui peuvent comprendre
		\item \expword{La capacité à construire de l'ADN à partir de zéro s'améliore de façon exponentielle}
		\item \expword{La capacité à construire de l'ADN à partir de zéro s'améliore plus rapidement que les ordinateurs} : Public moins expérimenté
	\end{itemize}
\end{minipage}
%
\begin{minipage}{0.39\textwidth}
	\hgraphpage{audience-experience-exp.png}
\end{minipage}

\end{frame}

\begin{frame}
\frametitle{Préparation : Audience}
\framesubtitle{Métier}

\begin{itemize}
	\item \optword{Présentation académique}
	\begin{itemize}
		\item audience : étudiants, enseignants, chercheurs, apprentis, etc.
		\item exemple : \expword{cours}, \expword{rapport d'un étudiant}, \expword{recherche}, \expword{séminaire}
		\item but : améliorer le savoir
		\item Plus formelle (exemple : utilisation de "nous" à la place de "je")
	\end{itemize}
	\item \optword{Présentation professionnelle}
	\begin{itemize}
		\item audience : client, investisseur, etc.
		\item exemple : \expword{produit}, \expword{idée}
		\item but : prendre de décisions
	\end{itemize}
\end{itemize}

\end{frame}

\subsection{Sujet de la présentation}

\begin{frame}
\frametitle{Préparation}
\framesubtitle{Sujet (Quoi?)}

\begin{itemize}
	\item Comprendre son sujet : l'objectif, le fond et comment atteindre ces objectifs
	\item Énumérer les concepts clés et les points à transmettre
	\item Commencez à réfléchir aux moyens d'illustrer les points clés
\end{itemize}

\end{frame}

%===================================================================================
\section{Contenu d'une présentation}
%===================================================================================

\subsection{Structure d'une présentation}

\begin{frame}
\frametitle{Contenu}
\framesubtitle{Structure}

\begin{itemize}
	\item \optword{Début} 
	\begin{itemize}
		\item expliquer le but de cette présentation
		\item Qui? Quoi? Où? Quand? Pourquoi?
	\end{itemize}
	\item \optword{Corps}
	\begin{itemize}
		\item expliquer le sujet
		\item Comment?
	\end{itemize}
	\item \optword{Conclusion}
	\begin{itemize}
		\item résumer le sujet
	\end{itemize}
\end{itemize}

\end{frame}

\begin{frame}
\frametitle{Contenu : Structure}
\framesubtitle{Début}

\begin{itemize}
	\item Le début doit être bien fait
	\item Il doit convaincre l'audience qu'on 
	\begin{itemize}
		\item ne va pas gaspiller son temps
		\item est bien organisé
		\item connait son audience
		\item maitrise son sujet
	\end{itemize}
\end{itemize}

\end{frame}

\begin{frame}
\frametitle{Contenu : Structure}
\framesubtitle{Début en détail}

\begin{itemize}
	\item Se présenter (première diapositive)
	\begin{itemize}
		\item Titre du sujet ; Nom et prénom ; Affiliation (école, université, laboratoire, entreprise, etc.) ; Date
	\end{itemize}

	\item Contexte du sujet (des fois, pas nécessaire)
	\begin{itemize}
		\item des informations qui aident l'audience à comprendre le thème du sujet
		\item c-à-d, une introduction à la problématique 
		\item exemple : \expword{décrire la façon de travail d'une bibliothèque }
	\end{itemize}

	\item Motivation
	\begin{itemize}
		\item pourquoi notre
		\item la problématique du sujet
		\item les objectifs du travail (présentation)
		\item exemple : \expword{automatiser la gestion des livres dans cette bibliothèque}
		\item Essayer de poser des questions pour stimuler la réflexion
	\end{itemize}

	\item Plan de la présentation
	\begin{itemize}
		\item Énumérer ce qu'on va présenter 
	\end{itemize}
\end{itemize}

\end{frame}

\begin{frame}
\frametitle{Contenu : Structure}
\framesubtitle{Corps}

\begin{itemize}
	\item Présenter des données et des faits
	\item Incorporer des citations d'experts
	\item Racontez ses expériences personnelles
	\item Être crédible
	\begin{itemize}
		\item être précis et exact avec les citations, les noms, les dates et les faits.
		\item fournir des références qui supportent les idées présentées
		\item insérer le contenu qu'on maitrise 
		\item utiliser un langage approprié à l'audience
	\end{itemize}
	\item Les diapositives doivent être homogènes 
	\item Il ne faut pas mettre trop d'informations sur une même diapositive
\end{itemize}

\end{frame}

\begin{frame}
\frametitle{Contenu : Structure}
\framesubtitle{Corps d'une présentation académique}

\begin{itemize}
	\item Travaux connexes
	\begin{itemize}
		\item Décrire des travaux similaires (les différentes méthodes pour répondre à la même problématique)
	\end{itemize}
	\item Méthodes
	\begin{itemize}
		\item Décrire son approche : les méthodes proposées 
		\item Illustrer les méthodes par des schémas
	\end{itemize}
	\item Synthèse
	\begin{itemize}
		\item Décrire les tests et les méthodes de validation 
		\item Essayer de présenter l'essentiel des résultats (pas le tout)
		\item Synthétiser : pourquoi la méthode donne ces résultats 
		\item Comparer son méthode avec des travaux similaires
	\end{itemize}
\end{itemize}

\end{frame}

\begin{frame}
\frametitle{Contenu : Structure}
\framesubtitle{Conclusion}

\begin{itemize}
	\item Reprendre la motivation de cette présentation
	\item Résumer les points principaux d'une façon simple
	\item Discuter les résultats
	\item Préciser comment les méthodes proposées ont réglées le problème
	\item Énumérer les limites et comment améliorer le travail
\end{itemize}

\end{frame}

\subsection{Esthétique d'une présentation}

\begin{frame}
\frametitle{Contenu}
\framesubtitle{Esthétique}

\begin{itemize}
	\item Les illustrations
	\item Les polices et texte
	\item Les couleurs
	\item Les animations 
\end{itemize}

\end{frame}


\begin{frame}
\frametitle{Contenu : Esthétique}
\framesubtitle{Illustrations}

\begin{minipage}{0.60\textwidth}
\begin{itemize}
	\item Enrichir la présentation par  des illustrations : images et tableaux
	\item L'audience a tendance d'apprendre plus en entendant et regardant au même temps
	\item Plus d'illustrations, moins de texte
	\item Il faut :
	\begin{itemize}
		\item fournir des titre aux illustrations
		\item s'assurer que les illustrations sont lisibles.
		\item citer l'origine des illustrations
	\end{itemize}
\end{itemize}
\end{minipage}
%
\begin{minipage}{0.39\textwidth}
	\begin{figure}
		\hgraphpage{illustrations.pdf}
		\caption{Pourcentage de rappel selon \cite{cipolla}}
	\end{figure}
\end{minipage}

\end{frame}

\begin{frame}
\frametitle{Contenu : Esthétique}
\framesubtitle{Polices et texte}

\begin{itemize}
	\item Utiliser des grandes tailles de police (plus de 18pt) afin d'améliorer la lisibilité 
	\item Utiliser des tailles différentes pour les titres et les sous titres 
	\item Utiliser une police standard comme \expword{Arial} ; n'utiliser pas des polices difficiles à lire
	\item Le texte dans une diapositive : 
	\begin{itemize}
		\item sous forme de listes à puce
		\item le minimum possible
	\end{itemize}
\end{itemize}

\end{frame}

\begin{frame}
\frametitle{Contenu : Esthétique}
\framesubtitle{Polices et texte : exemple à ne pas faire}

{\tiny 

Ceci est un texte trop petit et difficile à lire par l'audience. 
J'ai mis ce texte comme ça afin de vous torturer ; vous finissez par porter des lunettes. 
Même en les portant, vous ne pouvez pas lire ce texte à distance puisque j'ai rempli cette diapositive par du n'importe quoi. 
Ce texte est trop long, même moi j'avez du mal à l'écrire.

Ceci est un texte trop petit et difficile à lire par l'audience. 
J'ai mis ce texte comme ça afin de vous torturer ; vous finissez par porter des lunettes. 
Même en les portant, vous ne pouvez pas lire ce texte à distance puisque j'ai rempli cette diapositive par du n'importe quoi. 
Ce texte est trop long, même moi j'avez du mal à l'écrire.

Ceci est un texte trop petit et difficile à lire par l'audience. 
J'ai mis ce texte comme ça afin de vous torturer ; vous finissez par porter des lunettes. 
Même en les portant, vous ne pouvez pas lire ce texte à distance puisque j'ai rempli cette diapositive par du n'importe quoi. 
Ce texte est trop long, même moi j'avez du mal à l'écrire.

Ceci est un texte trop petit et difficile à lire par l'audience. 
J'ai mis ce texte comme ça afin de vous torturer ; vous finissez par porter des lunettes. 
Même en les portant, vous ne pouvez pas lire ce texte à distance puisque j'ai rempli cette diapositive par du n'importe quoi. 
Ce texte est trop long, même moi j'avez du mal à l'écrire.

Ceci est un texte trop petit et difficile à lire par l'audience. 
J'ai mis ce texte comme ça afin de vous torturer ; vous finissez par porter des lunettes. 
Même en les portant, vous ne pouvez pas lire ce texte à distance puisque j'ai rempli cette diapositive par du n'importe quoi. 
Ce texte est trop long, même moi j'avez du mal à l'écrire.

}

{\calligra Pour mieux vos torturer, j'ai écrit ce texte. En fin de compte, vous êtes fautifs pour atteindre ma présentation}

\end{frame}

\begin{frame}
\frametitle{Contenu : Esthétique}
\framesubtitle{Couleurs}

\begin{itemize}
	\item La couleur de police doit être opposée à celle de l'arrière plan
	\item Utilisez les couleurs pour renforcer la logique de votre structure
	\item Les couleurs doivent être compatibles (si on utilise le rouge pour les mots clés, on ne doit pas utiliser une autre couleur)
	\item Par exemple, dans cette présentation : 
	\begin{itemize}
		\item \expword{Ceci est un exemple}
		\item \keyword{Ceci est un mot clé}
		\item \optword{Ceci est un choix dans une liste} 
	\end{itemize}
	\item Il ne faut pas exagérer dans l'utilisation des couleurs 
\end{itemize}

\end{frame}

{
\setbeamercolor{background canvas}{bg=red}
\begin{frame}
\frametitle{Contenu : Esthétique}
\framesubtitle{Couleurs : exemple à ne pas faire}

\textcolor{yellow}{Pour détruire vos yeux, j'ai choisi ces belles couleurs. Du rouge comme fond pour inciter vos yeux et du jaune pour l'écriture afin que vous concentriez plus en la lisant. Une combinaison mortelle.}

\vfill

\textcolor{WildStrawberry}{J'ai même choisi une couleur presque identique à celle du fond pour les gens avec des yeux solides}

\end{frame}
}

\begin{frame}
\frametitle{Contenu : Esthétique}
\framesubtitle{Animations}

\begin{itemize}
	\item L'intérêt : 
	\begin{itemize}
		\item présenter d'une façon séquentielle
		\item éviter de dévoiler tout ce qu'on présente sur la même diapositive
		\item guider l'œil 
	\end{itemize}
		\item Il ne faut pas utiliser des : 
	\begin{itemize}
		\item différentes animations au même type de texte (par exemple, titre)
		\item animations complexes pour une présentation académique (\keyword{Keep it simple})
	\end{itemize}
\end{itemize}

\end{frame}

%===================================================================================
\section{Présentation orale}
%===================================================================================

\subsection{Avant la présentation}

\begin{frame}
\frametitle{Présentation orale}
\framesubtitle{Avant la présentation}

\begin{itemize}
	\item S'entrainer sur la présentation
	\begin{itemize}
		\item En face ces amis pour avoir des avis
		\item En utilisant un chronomètre pour mesurer le temps 
		\item Anticiper des questions probables et préparer des réponses 
	\end{itemize}
	\item Tester la présentation sur un data-show 
	\item Se préparer à sauter des diapositives si le temps ne sera pas suffisant
\end{itemize}

\end{frame}

\subsection{Durant la présentation}

\begin{frame}
\frametitle{Présentation orale}
\framesubtitle{Durant la présentation}

\begin{itemize}
	\item Énergie (pas au point d'être hyperactif)
	\begin{itemize}
		\item Utiliser un bon volume et intonation du son 
		\item Être passionné par son sujet et avoir confiance en soi
		\item Se déplacer. Ne rester pas au même endroit.
	\end{itemize}
	\item Attitude et langage du corps
	\begin{itemize}
		\item Contact visuel : Regarder l'audience dans les yeux
		\item Utilisez des expressions faciales positives telles que les sourires, les yeux expressifs et les regards empathiques
		\item Éviter des gestes comme les allers-retours, agiter les main, croiser les bras, etc. 
	\end{itemize}
\end{itemize}

\end{frame}

\begin{frame}
\frametitle{Présentation orale : Durant la présentation}
\framesubtitle{Un peu d'humour}

\begin{center}
	Si vous sentez que vous êtes entrain de perdre l'audience, utilisez cette technique

\vgraphpage[.7\textheight]{puss-in-boots.jpg}
\end{center}


\end{frame}

\subsection{Après la présentation}

\begin{frame}
\frametitle{Présentation orale}
\framesubtitle{Après la présentation}

\begin{itemize}
	\item Réfléchir avant de répondre
	\item Demander de clarification au cas où la question n'était pas claire 
	\item Il ne faut pas répondre à une question inattendue par une réponse vague et hésitante.
	Il faut mieux dire que c'est un point très intéressant sur lequel on devra enquêter, ou que c'est en dehors de son sujet 
	\item Garder les réponses aux questions brèves
\end{itemize}

\end{frame}

\subsection{Remarques à l'audience}

\begin{frame}
\frametitle{Présentation orale}
\framesubtitle{Remarques à l'audience}

\begin{itemize}
	\item L'audience est un critère de la réussite d'une présentation
	\item Participer à la conversation. On n'est pas entrain de regarder pas la télévision
	\item Penser à la façon dont le matériel présenté remet en question ce qu'on sait déjà
	\item Poser des questions
	\item Suivre la présentation sans faire de bruit
\end{itemize}

\end{frame}

%===================================================================================
\section{Microsoft PowerPoint}
%===================================================================================

\begin{frame}
\frametitle{Microsoft PowerPoint}
%\framesubtitle{}
\begin{itemize}
	\item logiciel de présentation
	\item fait partit à la suite bureautique \optword{Microsoft office}
	\item disponible sur Windows, Mac OS, iOS et Android
	\item le programme de présentation le plus utilisé dans le monde
\end{itemize}
\end{frame}

\begin{frame}
\frametitle{Microsoft PowerPoint}
\framesubtitle{Vue globale}

\begin{center}
	\vgraphpage{pp-general.png}
\end{center}

\end{frame}

\subsection{Préparation}

\begin{frame}
\frametitle{Microsoft PowerPoint}
\framesubtitle{Préparation}

\begin{itemize}
	\item une diapositive PowerPoint est une page vide 
	\item pour ajouter du texte, il faut ajouter une zone de texte
	\item les groupes \optword{police} et \optword{paragraphe} de l'onglet \optword{Accueil} ne peuvent être activés que lorsqu'on sélectionne une zone de texte 
	\item on peut insérer des diapositives avec des composants prédéfinies comme le titre et le contenu
\end{itemize}

\end{frame}

\begin{frame}
\frametitle{Microsoft PowerPoint : Préparation}
\framesubtitle{Disposition d'une diapositive}

\begin{minipage}{0.60\textwidth}
	\begin{itemize}
		\item Pour avoir une présentation homogène, on peut utiliser des dispositions pour chaque type de diapositives
		\item Quelques types : 
		\begin{itemize}
			\item \optword{Diapositive de titre} : la première diapositive de la présentation
			\item \optword{Titre de section} : pour introduire une nouvelle section. pae exemple, \expword{Tests et résultats}
			\item \optword{Titre et contenu} : une diapositive avec un titre et du contenu, comme celle-ci
		\end{itemize}
		\item Pour changer la disposition d'une diapostive déjà crée, cliquer sur \optword{Disposition} dans l'onglet \optword{Accueil}
	\end{itemize}
\end{minipage}
%
\begin{minipage}{0.39\textwidth}
	\hgraphpage{pp-nouv-diapo.png}
\end{minipage}

\end{frame}

\begin{frame}
\frametitle{Microsoft PowerPoint : Préparation}
\framesubtitle{Insertion des éléments}

\hgraphpage{pp-insertion.png}

\begin{itemize}
	\item Pour enrichir la diapositive, il faut insérer des éléments à partir de l'onglet \optword{Insertion}
	\item Les mêmes éléments que ceux de Microsoft Word
	\item Les présentations académiques doivent toujours être numérotées
	\item Pour afficher la numérotation, ainsi que d'autres informations, sur les diapositives : \optword{Entête/Pied}
\end{itemize}
\end{frame}

\begin{frame}
\frametitle{Microsoft PowerPoint : Préparation}
\framesubtitle{Thèmes des diapositives}

\hgraphpage{pp-creation.png}

\begin{minipage}{0.60\textwidth}
	\begin{itemize}
		\item On peut changer le thème de la présentation 
		\item On peut, aussi, modifier l'arrière plan d'une ou toutes les diapositives
		\item Selon le type du data-show, on peut définir la taille de la présentation : 
		\begin{itemize}
			\item standard TV (4:3)
			\item HDTV (16:9)
			\item Autres
		\end{itemize}
	\end{itemize}
\end{minipage}
%
\begin{minipage}{0.39\textwidth}
	\hgraphpage{pp-taille.png}
\end{minipage}

\end{frame}

\begin{frame}
\frametitle{Microsoft PowerPoint : Préparation}
\framesubtitle{Révision}

\hgraphpage{pp-revision.png}

\begin{itemize}
	\item Comme dans Microsoft Word, on peut réviser la présentation : 
	\begin{itemize}
		\item Vérification de l'orthographe  
		\item Commenter (commentaires de révision) et comparer deux présentations
	\end{itemize}
\end{itemize}

\end{frame}

\begin{frame}
\frametitle{Microsoft PowerPoint : Préparation}
\framesubtitle{Diapograma}

\hgraphpage{pp-diapograma.png}

\begin{minipage}{0.60\textwidth}
	\begin{itemize}
		\item 
	\end{itemize}
\end{minipage}
%
\begin{minipage}{0.39\textwidth}
	\hgraphpage{pp-diapograma-param.png}
\end{minipage}

\end{frame}


\subsection{Transitions et Animations}

\begin{frame}
\frametitle{Microsoft PowerPoint : Transitions et Animations}
\framesubtitle{Transitions}

\hgraphpage{pp-transition.png}

\begin{itemize}
	\item Lorsqu'on passe d'une diapositive vers une autre
	\item Appliquer une transition sur une ou toutes les diapositives
	\item On peut définir la durée de la transaction
	\item On peut ajouter du son avec la transaction (ne le fait pas dans les présentations des rapports)
	\item On peut passer vers une autre diapositive manuellement ou automatiquement après une durée
	\item éviter le passage automatique, puisque la présentation peut passer vers la diapositive suivante lorsqu'on est entrain de décrire la diapositive courante
\end{itemize}

\end{frame}


\begin{frame}[t]
\frametitle{Microsoft PowerPoint : Transitions et Animations}
\framesubtitle{Animations}

\hgraphpage{pp-animation.png}

\begin{minipage}{0.58\textwidth}
	\begin{itemize}
		\item Sélectionner un texte ou un objet
		\item Choisir une animation : 
		\begin{itemize}
			\item Entrée : afficher l'objet
			\item Accentuation : ni afficher ni cacher 
			\item Sortie : cacher l'objet
		\end{itemize}
		\item Démarrage des animations : 
		\begin{itemize}
			\item Au clic
			\item Avec la précédente
			\item Après la précédente 
		\end{itemize}
		\item Pour organiser les animations, afficher le volet des animations
	\end{itemize}
\end{minipage}
%
\begin{minipage}{0.20\textwidth}
	\hgraphpage{pp-animation-types.png}
\end{minipage}
\begin{minipage}{0.20\textwidth}
	\hgraphpage{pp-animation-volet.png}
\end{minipage}

\end{frame}

\subsection{Affichage}

\begin{frame}
\frametitle{Microsoft PowerPoint}
\framesubtitle{Affichage}

\hgraphpage{pp-affichage.png}

\begin{itemize}
	\item \optword{Normal} : le mode d'édition dans lequel on travaille le plus souvent pour créer les diapositives
	\item \optword{Mode Plan} : pour créer un plan de présentation. Il affiche uniquement le texte et pas les images ou autres éléments graphiques
	\item \optword{Trieuse de diapositives} : afficher toutes les diapositives sous forme de miniatures disposées horizontalement. Il est utile pour réorganiser les diapositives
	\item \optword{Page de notes} : diapositives avec volet \optword{Notes}. Les notes sont utilisées comme indications pendant la présentation.
	\item \optword{Mode Lecture} : Exécuter sans mode plein écran
\end{itemize}
\end{frame}


\begin{frame}[t]
\frametitle{Microsoft PowerPoint : Affichage}
\framesubtitle{Masque des diapositives}

\hgraphpage{pp-masque-diapo.png}

\begin{minipage}{0.80\textwidth}
\begin{itemize}
	\item Pour créer des dispositions personnalisées ou modifier celles existantes
	\item On peut ajouter des éléments à un type de diapositives
	\item Par exemple, \expword{le logo de l'université}
	\item Tous ce qu'on met sur le premier masque, s'applique sur les masques fils
\end{itemize}
\end{minipage}
%
\begin{minipage}{0.19\textwidth}
	\hgraphpage{pp-masque-diapo2.png}
\end{minipage}

\end{frame}

\begin{frame}[t]
\frametitle{Microsoft PowerPoint : Affichage}
\framesubtitle{Masque des pages}

\hgraphpage{pp-masque-doc.png}

\begin{minipage}{0.70\textwidth}
	\begin{itemize}
		\item Pour personnaliser comment PowerPoint va imprimer les diapositives
	\end{itemize}
\end{minipage}
%
\begin{minipage}{0.29\textwidth}
	\hgraphpage{pp-masque-doc2.png}
\end{minipage}

\end{frame}

\begin{frame}[t]
\frametitle{Microsoft PowerPoint : Affichage}
\framesubtitle{Masque des notes}

\hgraphpage{pp-masque-notes.png}

\begin{minipage}{0.70\textwidth}
	\begin{itemize}
		\item Pour personnaliser comment PowerPoint va imprimer les diapositives avec leurs notes
		\item Rappel : une note c'est le texte qui décrit la diapositive (le speech)
	\end{itemize}
\end{minipage}
%
\begin{minipage}{0.29\textwidth}
	\hgraphpage{pp-masque-notes2.png}
\end{minipage}

\end{frame}


\begin{frame}
\frametitle{Présentation}
\framesubtitle{Un peu d'humour}

\begin{center}
	\vgraphpage{humour.jpg}
\end{center}

\end{frame}

\insertbibliography{Bweb05}{*}


\end{document}

